\documentclass{article}
\usepackage{mathtools} 
\usepackage{fontspec}
\usepackage[UTF8]{ctex}
\usepackage{amsthm}
\usepackage{mdframed}
\usepackage{xcolor}
\usepackage{amssymb}
\usepackage{amsmath}
\usepackage{hyperref}
\usepackage{mathrsfs}



% 定义新的带灰色背景的说明环境 zremark
\newmdtheoremenv[
  backgroundcolor=gray!10,
  % 边框与背景一致，边框线会消失
  linecolor=gray!10
]{zremark}{注释}

% 通用矩阵命令: \flexmatrix{矩阵名}{元素符号}{行数}{列数}
\newcommand{\flexmatrix}[4]{
  \[
  #1 = \begin{pmatrix}
    #2_{11}     & #2_{12}     & \cdots & #2_{1#4}   \\
    #2_{21}     & #2_{22}     & \cdots & #2_{2#4}   \\
    \vdots      & \vdots      & \ddots & \vdots     \\
    #2_{#31}    & #2_{#32}    & \cdots & #2_{#3#4}
  \end{pmatrix}
  \]
}

% 简化版命令（默认矩阵名为A，元素符号为a）: \quickmatrix{行数}{列数}
\newcommand{\quickmatrix}[2]{\flexmatrix{A}{a}{#1}{#2}}


\begin{document}
\title{7.1}
\author{张志聪}
\maketitle

\section*{1}

todo

\section*{2}

$V_{\lambda_0}$特征子空间的维数为$k$，
于是可得$rank(A) = n - k$，由命题1.4可知，在$V$内存在一组基，
使得该组基下$\mathscr{A}$的矩阵成Jordan形矩阵，
因为$rank(A) = n - k$，所以$Jordan$形矩阵$J$中，只有$n - k$个1，
于是单个Jordan块($J_{ni}$)中$1$的个数不超过$n - k$个，
从而$J_{ni}$的阶数$\leq n - k + 1$，
于是有$J_{ni}^{n - k + 1} = 0$，
由分块矩阵乘法可知，$J^{n - k + 1} = 0$.

\section*{3}

tip:
判断一个矩阵是不是幂零矩阵，通常使用特征多项式是否为$f(\lambda) = \lambda^n$来判断.

以第一个矩阵为例。

我们有
\begin{align*}
  f(\lambda) = |\lambda E - A| = \lambda^3
\end{align*}

\begin{align*}
  A   & = \begin{bmatrix}
            0  & - 3 & 3  \\
            -2 & - 7 & 13 \\
            -1 & - 4 & 7
          \end{bmatrix} \\
  A^2 & = \begin{bmatrix}
            3 & 9 & -18 \\
            1 & 3 & -6  \\
            1 & 3 & -6
          \end{bmatrix} \\
  A^3 & = 0
\end{align*}

综上可得$\mathscr{A}$是循环幂零线性变换。
\section*{4}

按照题设，得
\begin{align*}
  \mathscr{A}^n X = \sum \limits_{i = 1}^n (-1)^i C_n^i A^{n - i} X A^i
\end{align*}
以上可对n进行归纳完成证明。

因为$A$是一个幂零矩阵，即存在整数$m$，使得$A^m = 0$.
由于$\max\{2m - i, i\} \geq m, \forall i = 0, 1, 2, \cdots, 2m$.
因此，对任意$X \in M_n(k), \mathscr{A}^{2m} X = 0$.

\section*{5}

 (i)由导数定义，对任意$f(x) \in K[x]_n$，有
\begin{align*}
  \mathscr{D}^n f = 0
\end{align*}
且对$f(x) = x^{n - 1}$，有
\begin{align*}
  \mathscr{D}^{n - 1} f = (n - 1)! \neq 0
\end{align*}
故$\mathscr{D}^{n - 1} \neq 0$。

综上，$\mathscr{D}$是循环幂零线性变换。

(ii)一组循环基：
\begin{align*}
  \mathscr{D}^{n - 1} x^{n - 1},
  \mathscr{D}^{n - 2} x^{n - 1},
  \cdots,
  \mathscr{D} x^{n - 1},
  x^{n - 1}
\end{align*}

\section*{6}

$\mathscr{A}$的全部不变子空间只有
\begin{align*}
  \{0\}, ker \mathscr{A}, ker \mathscr{A}^2, \cdots, ker \mathscr{A}^{n - 1}, V
\end{align*}
一共n + 1个。

因为$\mathscr{A}$是循环幂零线性变换，于是可得对于每个$k$，
\begin{align*}
  ker\, \mathscr{A}^k = span \{\mathscr{A}^{n - k} \alpha, \mathscr{A}^{n - k + 1} \alpha, \cdots, \mathscr{A}^{n - 1} \alpha\}
\end{align*}
所以
\begin{align*}
  \{0\} \subset ker \mathscr{A} \subset ker \mathscr{A}^{2} \subset \cdots \subset ker \mathscr{A}^{n - 1} \subset V
\end{align*}
另外，需要证明$\mathscr{A}$的任意非零不变子空间$M$，都属于以上子空间。

为了方便表述，令
\begin{align*}
  e_1 & = \mathscr{A}^{n-1} \alpha \\
  e_2 & = \mathscr{A}^{n-2} \alpha \\
  \vdots                           \\
  e_n & = \alpha
\end{align*}

$M$的一组基可以被$e_1, e_2, \cdots, e_n$表示，
取$M$任意基都可以表示成
\begin{align*}
  \epsilon_i = a_1 e_1 + a_2 e_2 + \cdots + a_n e_n, i = 1, 2, \cdots, n
\end{align*}
在该基下的最高非零坐标对应基向量$e_r$，此时$a_r \neq 0$。

由于$M$是不变子空间，于是我们有
\begin{align*}
  \mathscr{A}^{n - r} \epsilon_i & = a_r e_1
\end{align*}
可知$e_1 \in M$，因为
\begin{align*}
  \mathscr{A}^{n - r - 1} \epsilon_i
   & = a_{r - 1} e_1 + a_r e_2
\end{align*}
也是可得$e_2 \in M$，
类似地，可得$e_3, e_4, \cdots, e_r \in M$，故$ker \mathscr{A}^{r} \subset M$.

接下来，我们证明$M \subset ker \mathscr{A}^{r}$。

首先，证明$M$必须属于以上$n + 1$个不变子空间，
$M = V$结论是显然的；$M \neq V$，那么$M$是否属于某个
$ker \mathscr{A}, ker \mathscr{A}^{2}, \cdots, ker \mathscr{A}^{n - 1}$，
假设不属于，于是存在$\beta \in M$，
$\mathscr{A} \beta \neq 0, \mathscr{A}^{2} \beta \neq 0, \cdots, \mathscr{A}^{n - 1} \beta \neq 0$，
因为$\beta, \mathscr{A} \beta, \mathscr{A}^{2} \beta, \cdots, \mathscr{A}^{n - 1} \beta$线性无关，
所以$M$是n维空间，故$M = V$，与前提矛盾，所以$M$属于某个
$ker \mathscr{A}, ker \mathscr{A}^{2}, \cdots, ker \mathscr{A}^{n - 1}$.

假设$M \not \subset ker \mathscr{A}^{r}$，
于是可得$M \subset ker \mathscr{A}^{s}, s > r$（$M = V$表述依然成立），
所以$e_s \in M$，这与$M$在该基下的最高非零坐标对应基向量$e_r$矛盾，
假设不成立，即$M \subset ker \mathscr{A}^{r}$.

综上$M = ker \mathscr{A}^{r}$，命题得证。

\section*{7}

$\mathscr{A}$是幂零线性变换，故$\mathscr{A}$只有特征值0，所以
$\alpha, \beta$属于$V_0$特征子空间，
由题设可知$V_0$的维数$\geq 2$，不妨设为$k$，
由习题2可知，$\mathscr{A}^{n - k + 1} = 0$，
由于$n - k + 1 < n - 1$，由循环幂零线性变换定义，
$\mathscr{A}$不是循环幂零线性变换。

\section*{8}

由$\mathscr{A} \mathscr{B} = \mathscr{B} \mathscr{A}$可知，

接下来证明，若干$\mathscr{A}, \mathscr{B}$的构成的乘积，与排列顺序无关。

假设有$m$个$\mathscr{A}$，$n$个$\mathscr{B}$，$t = m + n$，对$t$进行归纳。

$t = 1$,显然成立，因为要么只有一个$\mathscr{A}$，要么只有一个$\mathscr{B}$;

归纳假设$< t$时，命题成立。

考虑一个有$m$个$\mathscr{A}$，$n$个$\mathscr{B}$组成的，长度为$t$的乘积：
\begin{align*}
  P = x_1 x_2 \cdots x_t
\end{align*}
其中$x_i$是$\mathscr{A}$或$\mathscr{B}$.

如果$x_1 = \mathscr{A}$，则
\begin{align*}
  P = \mathscr{A} \cdot (x_2 x_3 \cdots x_t)
\end{align*}
于是$(x_2 x_3 \cdots x_t)$中的总数为$t - 1$项，其中有$m - 1$个$\mathscr{A}$，$n$个$\mathscr{B}$，
所以由归纳假设可知
\begin{align*}
  P & = \mathscr{A} \cdot \mathscr{A}^{m - 1} \mathscr{B}^n \\
    & = \mathscr{A}^m \mathscr{B}^n
\end{align*}
同理$x_1 = \mathscr{B}$时，
\begin{align*}
  P & = \mathscr{B} \cdot (x_2 x_3 \cdots x_t)              \\
    & = \mathscr{B} \cdot \mathscr{A}^m \mathscr{B}^{n - 1}
\end{align*}
由$\mathscr{A} \mathscr{B} = \mathscr{B} \mathscr{A}$可知，
$\mathscr{B} \cdot \mathscr{A}^m = \mathscr{A}^m \cdot \mathscr{B}$，所以
\begin{align*}
  P & = \mathscr{A}^m \cdot \mathscr{B} \mathscr{B}^{n - 1} \\
    & = \mathscr{A}^m \mathscr{B}^{n}
\end{align*}
因此，无论第一个因子是$\mathscr{A}$还是$\mathscr{B}$，结果都是$\mathscr{A}^m \mathscr{B}^{n}$。

归纳完成。

于是，我们有
\begin{align*}
  (\mathscr{A} + \mathscr{B})^n
   & = \sum \limits_{i = 0}^n C_n^i \mathscr{A}^{n - i} \mathscr{B}^i
\end{align*}

由题意可知，存在正整数$k_1, k_2$使得$\mathscr{A} = 0, \mathscr{B} = 0$，
取$l = 2 \max\{k_1, k_2\}$，于是我们有
\begin{align*}
  max\{2l - i, i\} \geq l, i = 0, 1, 2, \cdots, 2l
\end{align*}
于是
\begin{align*}
  (\mathscr{A} + \mathscr{B})^{2l} = 0
\end{align*}
所以，$\mathscr{A} + \mathscr{B}$是幂零线性变换。

\section*{9}

 (i)
假设$k \mathscr{E} - \mathscr{A}$不可逆，即$k \mathscr{E} - \mathscr{A}$存在非零解$\alpha$，
即
\begin{align*}
  (k \mathscr{E} - \mathscr{A}) \alpha & = 0                   \\
  k \alpha                             & =  \mathscr{A} \alpha
\end{align*}
于是可知$\mathscr{A}$存在特征值$k \neq 0$且$\alpha$是特征向量，
这与$\mathscr{A}$是幂零线性变换只有特征值0矛盾。

(ii)

取$\mathscr{A}$的一组基，设
$k \mathscr{E} - \mathscr{A}$在这组基下的矩阵为$A$，
由(i)可知矩阵$A$有可逆矩阵$A^{-1}$，由矩阵和线性变换一一对应的关系即可得到目标矩阵。

\section*{10}

有命题1.2可知
\begin{align*}
  V = I(\alpha_1) \oplus I(\alpha_2) \oplus \cdots \oplus I(\alpha_s)
\end{align*}
设$J_i$的阶数为$n_i$，于是我们有
\begin{align*}
  \mathscr{A}^{n_i} \alpha_i = 0
\end{align*}
即
\begin{align*}
  \mathscr{A} (\mathscr{A}^{n_i - 1} \alpha_i) = 0
\end{align*}
所以$\mathscr{A}^{n_i - 1} \alpha_i \in ker \mathscr{A} = V_{\lambda_0}, i = 1, 2, \cdots, s$.

接下来证明，$\mathscr{A}^{n_i - 1}, i = 1, 2, \cdots, s$是$V_{\lambda_0}$的一组基。

由$V = I(\alpha_1) \oplus I(\alpha_2) \oplus \cdots \oplus I(\alpha_s)$可知$\mathscr{A}^{n_i - 1}, i = 1, 2, \cdots, s$
线性无关。

假设$V_{\lambda_0}$的维数大于$k$，
即$I(\alpha_i), i = 1, 2, \cdots, s$向量是$V_{\lambda_0}$基向量，
不妨设$\beta \in I(\alpha_j), 1 \leq j \leq s$，
$\beta$可表示为
\begin{align*}
  \beta = a_1 \alpha_j + a_2 \mathscr{A} \alpha_j + \cdots + a_{n_j} \mathscr{A}^{n_j - 1} \alpha_j
\end{align*}
其中$a_i, i = 1, 2, \cdots, n_j - 1$不全为0（注意下标范围为什么这样取，因为全零，这$\beta = a_{nj}\mathscr{A}^{n_j - 1} \alpha_j$
已经在$\mathscr{A}^{n_i - 1}, i = 1, 2, \cdots, s$中了）。
因为$\beta \in V_{\lambda_0}$，所以
\begin{align*}
  \mathscr{A} \beta                                                                                           & = 0  \\
  \mathscr{A} (a_1 \alpha_j + a_2 \mathscr{A} \alpha_j + \cdots + a_{n_j} \mathscr{A}^{n_j - 1} \alpha_j)     & = 0  \\
  a_1 \mathscr{A} \alpha_j + a_2 \mathscr{A}^2 \alpha_j + \cdots + a_{n_j} \mathscr{A}^{n_j} \alpha_j         & = 0  \\
  a_1 \mathscr{A} \alpha_j + a_2 \mathscr{A}^2 \alpha_j + \cdots + a_{n_j - 1} \mathscr{A}^{n_j} \alpha_j + 0 & =  0 \\
  a_1 \mathscr{A} \alpha_j + a_2 \mathscr{A}^2 \alpha_j + \cdots + a_{n_j - 1} \mathscr{A}^{n_j} \alpha_j     & = 0
\end{align*}
这和$\mathscr{A} \alpha_j, \mathscr{A}^2, \cdots, \mathscr{A}^{n_j - 1}$线性无关矛盾，
故假设不成立。

综上，$V_{\lambda_0}$的维数等于$k$.

\section*{11}

todo

\section*{12}

以第一个矩阵为例。

由习题3可知，$\mathscr{A}$是循环幂零线性变换，
于是存在$\alpha \in V$，使得
\begin{align*}
  \alpha, \mathscr{A} \alpha, \mathscr{A}^2 \alpha
\end{align*}
组成一组基。这里需要注意$\alpha \not \in ker \mathscr{A}$，否则
$\mathscr{A} \alpha = 0$.我们通过解
\begin{align*}
  A x = 0
\end{align*}
这个线性方程组，解得
\begin{align*}
  ker  \mathscr{A} = span\left\{\begin{bmatrix}1 \\ 1 \\ 1\end{bmatrix}\right\}
\end{align*}
可得$dim \, ker \mathscr{A} = 1$，由习题10也可知，只有一个Jordan块。
于是可取
\begin{align*}
  \alpha               & = (\epsilon_1, \epsilon_2, \epsilon_3) \begin{bmatrix}
                                                                  0 \\
                                                                  0 \\
                                                                  1 \\
                                                                \end{bmatrix}   \\
  \mathscr{A} \alpha   & = (\epsilon_1, \epsilon_2, \epsilon_3) A \begin{bmatrix}
                                                                    0 \\
                                                                    0 \\
                                                                    1 \\
                                                                  \end{bmatrix}
  = (\epsilon_1, \epsilon_2, \epsilon_3) \begin{bmatrix}
                                           3  \\
                                           13 \\
                                           7
                                         \end{bmatrix}                          \\
  \mathscr{A}^2 \alpha & = (\epsilon_1, \epsilon_2, \epsilon_3) A \begin{bmatrix}
                                                                    3  \\
                                                                    13 \\
                                                                    7
                                                                  \end{bmatrix}
  = (\epsilon_1, \epsilon_2, \epsilon_3) \begin{bmatrix}
                                           - 18 \\
                                           -6   \\
                                           -6
                                         \end{bmatrix}                          \\
  \mathscr{A}^3 \alpha & = (\epsilon_1, \epsilon_2, \epsilon_3) \begin{bmatrix}
                                                                  - 18 \\
                                                                  -6   \\
                                                                  -6
                                                                \end{bmatrix}
  = 0
\end{align*}
所以以$\mathscr{A}^2 \alpha, \mathscr{A} \alpha, \alpha$为基，$\mathscr{A}$在此组基下的矩阵为
\begin{align*}
  J = \begin{bmatrix}
        0 & 1 & 0 \\
        0 & 0 & 1 \\
        0 & 0 & 0
      \end{bmatrix}
\end{align*}

\section*{13}

\begin{itemize}
  \item (1)

        我们有
        \begin{align*}
          \mathscr{T}^2(X)
           & = \mathscr{T}(\mathscr{T}(X)) \\
           & = \mathscr{T}(T^T X T)        \\
           & = T^T (T^T X T) T             \\
           & = (T^T)^2 X T^2
        \end{align*}
        容易得到$\mathscr{T}^k(X) = (T^T)^k X T^k$.
        有题设$T$是幂零矩阵可知，存在$k_1$使得$T^{k_1} = 0$，故
        \begin{align*}
          \mathscr{T}^{k_1}(X) = (T^T)^{k_1} X T^{k_1} & = 0
        \end{align*}
  \item (2)

        我们有
        \begin{align*}
          T^2 = 0
        \end{align*}
        所以
        \begin{align*}
          \mathscr{T}^2(X) = (T^T)^2 X T^2 = 0
        \end{align*}
        故$\mathscr{T}$是循环幂零线性变换。

        可选
        \begin{align*}
          X_0               & = \begin{bmatrix}
                                  1 & 0 \\
                                  0 & 1
                                \end{bmatrix} \\
          \mathscr{T} (X_0) & =
          \begin{bmatrix}
            0 & 0 \\
            1 & 0
          \end{bmatrix}
          \begin{bmatrix}
            1 & 0 \\
            0 & 1
          \end{bmatrix}
          \begin{bmatrix}
            0 & 1 \\
            0 & 0
          \end{bmatrix}
          = \begin{bmatrix}
              0 & 0 \\
              0 & 1
            \end{bmatrix}
        \end{align*}
        $X_0 \neq 0$，且$\mathscr{T}(X_0) \neq 0$，以$\mathscr{T}(X_0), X_0$作为一组基，
        $\mathscr{T}$在此组基下的矩阵为
        \begin{align*}
          \begin{bmatrix}
            0 & 1 \\
            0 & 0
          \end{bmatrix}
        \end{align*}
  \item (3)

        与(1)可知
        \begin{align*}
          \mathscr{T}^k(X) = (T^T)^k X T^k
        \end{align*}
        由题设$\mathscr{T}$是幂零线性变换，所以存在$k_1 \in \mathbb{N}$使得
        \begin{align*}
          \mathscr{T}^{k_1}(X) = (T^T)^{k_1} X T^{k_1} = 0, \forall X \in M_n(K)
        \end{align*}
        接下来证明$T^{k_1} = 0$，反证法，假设$T^{k_1} \neq 0$，那么$(T^T)^{k_1} \neq 0$，不妨设
        $(T^T)^{k_1}$的$i$列不为零，$T^{k_1}$的$j$行不为零，
        由$X$的任意性，取$X = E_{ij}$，即只在$i$行$j$列的位置上矩阵为1，其他位置为0。于是
        \begin{align*}
          (T^T)^{k_1} X T^{k_1} 
          & = (T^T)^{k_1} E_{ij} T^{k_1} \\
          & \neq 0
        \end{align*}
        出现矛盾，故$T^{k_1} = 0$。

        综上，$T$是幂零矩阵。
\end{itemize}

\section*{14}

反证法，假设存在大于$k$阶的Jordan块$J_i$，
即在某组基$\epsilon_1, \epsilon_2, \cdots, \epsilon_n$下，
$\mathscr{A}$是Jordan形矩阵$J$，
按照矩阵乘法和存在大于$k$阶的Jordan块，我们有
\begin{align*}
  J^k \neq 0
\end{align*}
这与$\mathscr{A}^k = 0$矛盾。

\end{document}
